\section{Objetivos}

O presente trabalho tem como objetivo principal investigar e propor melhorias nas formulações matemáticas existentes para o problema de compra mínima, com foco no aprimoramento de sua consistência estrutural e eficiência computacional. Em particular, busca-se desenvolver variações das formulações apresentadas que resultem em modelos mais fortes do ponto de vista de relaxação linear e, consequentemente, mais eficientes quando resolvidos por solvers de otimização modernos.

Nesse contexto, pretende-se analisar criticamente as formulações DIB, DD, DIV e DDI, identificando possíveis fragilidades, redundâncias ou oportunidades de fortalecimento por meio da introdução de novas restrições, reformulações ou técnicas de linearização. O objetivo é obter modelos que apresentem melhor desempenho prático, reduzindo o tempo de resolução e melhorando a qualidade dos limites obtidos por métodos exatos, especialmente quando utilizados em ferramentas como o solver Gurobi~\cite{gurobi}.

Adicionalmente, este trabalho visa abordar o desafio da escalabilidade, propondo abordagens heurísticas capazes de lidar com instâncias de grande porte, envolvendo mais de mil itens ou compradores. Para isso, serão desenvolvidas e avaliadas heurísticas baseadas no Algoritmo Genético de Chaves Aleatórias Enviésadas e na abordagem Otimizador de Chaves Aleatórias, explorando sua capacidade de produzir soluções de boa qualidade em tempos computacionais reduzidos.

Por fim, espera-se que a combinação entre melhorias nas formulações exatas e o desenvolvimento de métodos heurísticos contribua para ampliar a aplicabilidade prática do problema estudado, permitindo a resolução eficiente de instâncias que são desafiadoras para abordagens puramente exatas.